\chapter{全程短剧：从字节流到刀口}
\label{ch:walk}

本章用\textbf{同一段迷你字节流}分别演出 Hom / Chain / Native（TopoPH 类似 Hom+滤网）。\\
数字是教具，不是生产 cut。

\section{道具表}

字节流（十六进制示意）：
\begin{center}
\texttt{41 41 42 43 43 44 41 42 42 45 46 41}
\end{center}
对应颜色键（玩具）：A A B C C D A B B E F A

\begin{figure}[htbp]
\centering
\begin{tikzpicture}[font=\scriptsize\ttfamily]
  \foreach \i/\b/\c in {
    0/41/red!25,1/41/red!25,2/42/blue!25,3/43/green!30,
    4/43/green!30,5/44/yellow!35,6/41/red!25,7/42/blue!25,
    8/42/blue!25,9/45/orange!35,10/46/cyan!30,11/41/red!25} {
    \node[bytecell, fill=\c] at (\i*0.92,0) {\b};
    \node[font=\tiny, below] at (\i*0.92,-0.55) {\i};
  }
\end{tikzpicture}
\caption{教具字节流与颜色。}
\label{fig:toy-stream}
\end{figure}

\section{第一幕：Hom}

假设窗宽 5。走到位置 5 时窗为 \texttt{A A B C C}。\\
\textbf{假装}指纹叮了；异色边=2。\\
下一步到位置 6：新窗 \texttt{A B C C D}（若按进 D），边数变 $\Rightarrow$ 若叮了就切。

\begin{figure}[htbp]
\centering
\begin{tikzpicture}[font=\small]
  \node[softbox=Gen3Color] (1) {滚到 pos=6};
  \node[softbox=OkGreen, right=0.5cm of 1] (2) {叮！};
  \node[softbox=OkGreen, right=0.5cm of 2] (3) {边 2$\to$3};
  \node[softbox=CutRed, right=0.5cm of 3] (4) {咔嚓};
  \draw[-{Stealth}, thick] (1)--(2)--(3)--(4);
\end{tikzpicture}
\caption{Hom 短剧结局。}
\label{fig:act-hom}
\end{figure}

\section{第二幕：Chain}

同一流，密码锁状态在玩具里用 3 位表示。\\
走到某步低 3 位全 0（对齐）——才打开珠串检查。\\
对齐之间即使珠串大变也\textbf{不理}。

\begin{figure}[htbp]
\centering
\begin{tikzpicture}[font=\footnotesize]
  \foreach \i in {0,...,11} {
    \node[bytecell] (b\i) at (\i*0.85,0) {};
  }
  \foreach \i in {0,1,2,4,5,6,7,8,10,11} {
    \node[font=\tiny, gray, above] at (\i*0.85,0.55) {只拧};
  }
  \node[font=\tiny, OkGreen, above] at (3*0.85,0.55) {对齐};
  \node[font=\tiny, OkGreen, above] at (9*0.85,0.55) {对齐};
  \draw[cutmark] (9*0.85,-0.7)--(9*0.85,0.35);
  \node[font=\footnotesize, CutRed, below] at (9*0.85,-0.85) {对齐且骨架/环通过};
\end{tikzpicture}
\caption{Chain：时间大部分花在``拧锁''上。}
\label{fig:act-chain}
\end{figure}

\section{第三幕：Native}

stride=4。只在位置 3,7,11 钉星（玩具假设 $Y$ 刚好落格）。\\
星上的 Flip：0, 0, 1 —— 在第三颗星触发切，然后清空。

\begin{figure}[htbp]
\centering
\begin{tikzpicture}[font=\small]
  \draw[thick] (0,0)--(10,0);
  \foreach \x/\lab in {2.5/星1, 5.5/星2, 8.5/星3} {
    \fill[Gen51Color] (\x,0) circle (0.12);
    \node[above, font=\tiny] at (\x,0.15) {\lab};
  }
  \node[below, font=\tiny] at (2.5,-0.35) {Flip=0};
  \node[below, font=\tiny] at (5.5,-0.35) {0};
  \node[below, font=\tiny, CutRed] at (8.5,-0.35) {1 $\Rightarrow$ 切};
  \draw[cutmark] (8.5,0.4)--(8.5,-0.15);
\end{tikzpicture}
\caption{Native 短剧：事件在第三颗星响起。}
\label{fig:act-native}
\end{figure}

\section{插曲：挪 1 字节会发生什么}

把流整体右移 1 字节。尺子切会全盘错位。\\
Native / CDC 希望：裂缝粘在相同内容图案上，刀口``集体挪一下''后仍大面积重合（reuse）。

\begin{figure}[htbp]
\centering
\begin{tikzpicture}[font=\small]
  \draw[thick] (0,1)--(9,1);
  \draw[thick] (0,0)--(9,0);
  \foreach \x in {2,5,7.5} {
    \draw[OkGreen, very thick] (\x,0.85)--(\x,1.15);
    \draw[OkGreen, very thick] (\x+0.15,-0.15)--(\x+0.15,0.15);
  }
  \node[right, font=\footnotesize] at (9.1,1) {原流刀口};
  \node[right, font=\footnotesize] at (9.1,0) {$+1$B 后};
  \node[font=\footnotesize, gray] at (4.5,-0.7) {绿线对齐 = reuse 直觉};
\end{tikzpicture}
\caption{1-byte reuse：好 CDC 的``良心指标''。}
\label{fig:act-reuse}
\end{figure}

\section{导演备注}

\begin{itemize}[nosep]
  \item 真实系统有 min/max、hard cut、流式 feed、标定；
  \item 玩具省略了这些，只突出\textbf{判决内核}；
  \item 看懂短剧后再去正式手册翻公式，会轻松很多。
\end{itemize}
